\documentclass[12pt]{article}
\usepackage{amsmath, amssymb}
\usepackage[margin=1in]{geometry}
\setlength{\parskip}{6pt}
\title{QUBO Formulation: Recommender-System Feature Selection Problem}
\date{}
\begin{document}
\maketitle

\subsection*{Recommender-System Feature Selection Problem}

\paragraph{Problem Statement.}
Given a set of $N$ candidate features for a recommender model, each feature $i \in \{0, 1, \dots, N-1\}$ is associated with a utility score $u_i \geq 0$. Furthermore, each unordered pair of distinct features $(i,j)$ has an associated redundancy or complexity penalty $r_{ij} \geq 0$. The task is to select exactly $K$ features such that the total utility is maximized while the total redundancy penalty is minimized.

\paragraph{Decision Variables.}
For each candidate feature $i \in \{0, 1, \dots, N-1\}$, define a binary decision variable $x_i \in \{0, 1\}$, where $x_i = 1$ indicates that feature $i$ is selected, and $x_i = 0$ indicates it is not selected.

\paragraph{QUBO Derivation.}
Step 1 --- The original maximization objective is converted to a minimization problem by negating it:
\[
\min_{\mathbf{x} \in \{0,1\}^N} \left( -\sum_{i=0}^{N-1} u_i x_i + \sum_{0 \leq i < j \leq N-1} r_{ij} x_i x_j \right).
\]

Step 2 --- The cardinality constraint requiring exactly $K$ features to be selected is expressed algebraically as:
\[
\sum_{i=0}^{N-1} x_i = K.
\]

Step 3 --- To incorporate this equality constraint into the objective, we construct a quadratic penalty term with weight $P > 0$:
\[
\mathcal{P}(\mathbf{x}) = P \left( \sum_{i=0}^{N-1} x_i - K \right)^2.
\]
Expanding the square and applying the idempotent property $x_i^2 = x_i$ for binary variables yields:
\begin{align}
\left( \sum_{i=0}^{N-1} x_i - K \right)^2 &= \sum_{i=0}^{N-1} x_i^2 + \sum_{i \neq j} x_i x_j - 2K \sum_{i=0}^{N-1} x_i + K^2 \\
&= \sum_{i=0}^{N-1} x_i + 2\sum_{0 \leq i < j \leq N-1} x_i x_j - 2K \sum_{i=0}^{N-1} x_i + K^2 \\
&= (1 - 2K) \sum_{i=0}^{N-1} x_i + 2 \sum_{0 \leq i < j \leq N-1} x_i x_j + K^2.
\end{align}
Multiplying by $P$ gives the general penalty contribution:
\[
\mathcal{P}(\mathbf{x}) = P(1 - 2K) \sum_{i=0}^{N-1} x_i + 2P \sum_{0 \leq i < j \leq N-1} x_i x_j + P K^2.
\]

Step 4 --- Combining the negated objective and the penalty term yields the single QUBO Hamiltonian $H(\mathbf{x})$:
\[
H(\mathbf{x}) = \sum_{i=0}^{N-1} \left[ P(1 - 2K) - u_i \right] x_i + \sum_{0 \leq i < j \leq N-1} \left[ 2P + r_{ij} \right] x_i x_j + P K^2.
\]

Step 5 --- Reading off the coefficients from the standard QUBO form $H(\mathbf{x}) = \sum_{i} Q_{ii} x_i + \sum_{i < j} Q_{ij} x_i x_j + c$, the Q-matrix entries and offset are given in closed form as:
\begin{align}
Q_{ii} &= P(1 - 2K) - u_i, \quad \forall i \in \{0, \dots, N-1\}, \\
Q_{ij} &= 2P + r_{ij}, \quad \forall 0 \leq i < j \leq N-1, \\
Q_{ji} &= Q_{ij}, \quad \forall 0 \leq i < j \leq N-1, \\
c &= P K^2.
\end{align}

\paragraph{Penalty Weight Justification.}
The penalty weight $P$ must be chosen sufficiently large to ensure that any violation of the cardinality constraint incurs a higher energy cost than the maximum possible improvement in the objective function. The maximum utility gain achievable by adding or removing a single feature is bounded by $\max_i u_i$. Therefore, selecting $P > \max_i u_i$ guarantees that any deviation from the constraint $\sum_i x_i = K$ increases the total energy by at least $P$, which strictly dominates any potential reduction in the objective term. This ensures that the optimizer will never prefer an infeasible solution over a feasible one.

\paragraph{Correctness Argument.}
Minimizing $H(\mathbf{x})$ over $\{0,1\}^N$ recovers the optimal feature set because (i) for any assignment satisfying $\sum_i x_i = K$, the penalty term vanishes and $H(\mathbf{x})$ reduces exactly to the negated original objective; thus, minimizing $H$ over feasible points is equivalent to maximizing the original utility-minus-redundancy score. (ii) For any infeasible assignment, the penalty term is strictly positive and, by the choice of $P$, exceeds the maximum possible objective value achievable by any feasible assignment. Consequently, the global minimum of $H$ must occur at a feasible point that optimizes the original objective, guaranteeing that the QUBO formulation correctly encodes the combinatorial problem.

\paragraph{Illustrative Example.}
Consider a small instance with $N=3$ candidate features and a requirement to select exactly $K=2$ features. Let the utility scores be $u_0 = 5$, $u_1 = 3$, and $u_2 = 4$. Let the redundancy penalties be $r_{01} = 1$, $r_{02} = 2$, and $r_{12} = 1$. Following the justification, we choose $P = 6$ (since $\max_i u_i = 5$). 

The concrete objective to minimize is $\min -5x_0 - 3x_1 - 4x_2 + 1x_0x_1 + 2x_0x_2 + 1x_1x_2$. The concrete penalty term is $6(x_0 + x_1 + x_2 - 2)^2$, which expands to $-18x_0 - 18x_1 - 18x_2 + 12x_0x_1 + 12x_0x_2 + 12x_1x_2 + 24$. Substituting these values into the general Q-matrix formulas yields:
\begin{align}
Q_{00} &= 6(1 - 4) - 5 = -23, & Q_{11} &= 6(1 - 4) - 3 = -21, & Q_{22} &= 6(1 - 4) - 4 = -22, \\
Q_{01} &= 2(6) + 1 = 13, & Q_{02} &= 2(6) + 2 = 14, & Q_{12} &= 2(6) + 1 = 13, \\
c &= 6 \cdot 2^2 = 24.
\end{align}
The resulting concrete Hamiltonian is:
\[
H(x_0, x_1, x_2) = -23x_0 - 21x_1 - 22x_2 + 13x_0x_1 + 14x_0x_2 + 13x_1x_2 + 24.
\]
Enumerating the three feasible assignments (exactly two variables equal to 1):
\begin{itemize}
\item $\mathbf{x} = (1, 1, 0)$: $H = -23 - 21 + 13 + 24 = -7$ (Objective value: $5 + 3 - 1 = 7$).
\item $\mathbf{x} = (1, 0, 1)$: $H = -23 - 22 + 14 + 24 = -7$ (Objective value: $5 + 4 - 2 = 7$).
\item $\mathbf{x} = (0, 1, 1)$: $H = -21 - 22 + 13 + 24 = -6$ (Objective value: $3 + 4 - 1 = 6$).
\end{itemize}
The global minimum energy is $-7$, achieved by selecting features $\{0, 1\}$ or $\{0, 2\}$, which correspond to the maximum total score of $7$.

\end{document}